\section{Pascal-P e o p-code}
\label{sec:hist-pcode}

\paragraph{A ideia.}
O livro de \textcite{livro:pemberton1982}, que comenta linha a linha o
compilador Pascal-P4, resume o desenho logo na introdução: em vez de
gerar código para uma máquina real, o compilador gera o que veio a ser
chamado de \termo{p-code}{\emph{Pseudo-code}: instruções de uma máquina
hipotética, a P-machine, executadas por um interpretador.}, \enquote{for a
hypothetical stack-based computer that is in many ways ideal for Pascal
compilation}. O compilador e o interpretador do p-code são escritos em
Pascal.

\paragraph{Por que isso portava.}
Para levar o sistema a uma máquina nova, o livro descreve caminhos como
este: compilar o compilador Pascal-P com qualquer compilador Pascal já
existente, usá-lo para traduzir o próprio compilador para p-code e, daí
em diante, só traduzir à mão o \emph{interpretador} para a máquina nova.
O livro diz que o Pascal-P foi a base de muitos sistemas Pascal,
\enquote{including the well-known UCSD system}.

\paragraph{A máquina.}
O cap.~10 descreve a P-machine: um processador e uma memória, em que
muitas operações usam a pilha, chamadas de procedimento empilham
parâmetros e endereço de retorno, e um conjunto pequeno de registradores
(contador de programa, ponteiro de pilha e outros) controla a execução.

\paragraph{O que isso explica na JVM.}
\begin{itemize}
\item \emph{Bytecode de pilha.} A JVM também é uma máquina de pilha:
  \texttt{iload}, \texttt{iadd}, \texttt{ireturn} empilham e desempilham
  operandos, como a seção~\ref{sec:javap-basico} mostrou no
  \texttt{javap}. A troca, segundo \textcite{paper:vm-showdown}: código
  de pilha é menor, mas precisa de mais instruções para a mesma conta
  (seção~\ref{sec:hist-dalvik}).
\item \emph{Portar a VM, não o programa.} O \enquote{escreva uma vez, rode
  em qualquer lugar} do Java é a mesma aposta: o \texttt{.class} é o
  mesmo, e o trabalho de cada plataforma é ter uma JVM.
\item \emph{O custo.} Interpretar tem um preço por instrução:
  \textcite{paper:vm-showdown} descrevem cada despacho de instrução como
  um desvio indireto caro e, em geral, imprevisível. A JVM paga isso só
  no começo e depois compila; as seções~\ref{sec:hist-smalltalk}
  e~\ref{sec:hist-self} mostram de onde veio essa saída.
\end{itemize}
